\begin{textsample}{NEG Dim 2 – global\_north\_2023\_12 – Score 1.00 – global\_north\_2023\_12/103924868.txt}  \label{ex:f2_neg_205}
Pro-Palestine rally gives politicians fuel for anti-protest agenda A protest against the families of Israeli hostages has been turned into a political stunt, providing our leaders with a reason to call for increased anti-protest laws.
Tom Tanuki reports.
THERE ' S BEEN an outcry since Thursday over a group of Melbourne activists who, we ' ve been told, protested in Crowne Plaza against the families of Israeli hostages held since 7 October.
This outcry has been led by figures as senior as the Prime Minister and Opposition Leader themselves."What we saw last night in Melbourne at a hotel in Docklands goes beyond the right of people to peacefully protest in our democratic country.
There is no excuse, no circumstances where people should organise a demonstration against grieving families."Local ultra-Zionists like Menachem Vorchheimer have also been urging Victoria Police to take action, sharing personal details of activists from the Whistleblowers, Activists and Communities Alliance ( WACA ) -- who along with a coalition of activist groups undertook the action -- to encourage their arrest. ( Or whatever other action Menachem @ @ @ @ @ @ @ @ @ @. ) Incredible that these activists would stoop this low.
Targeting the families of hostages.
Can you believe it?
I didn't genuinely expect everyone to believe it and I initially commented as such online.
But I understand if you believed it.
There ' s footage of activists in the lobby and we ' re told the families were there.
So therefore, it seems, the activists undertook an action targeting the families.
So let ' s be clear.
Nobody created an action around protesting families of Israeli hostages.
This would not be a good idea from any serious person ' s perspective.
It would, to quote Jeremy Leibler, be"vile"to do it ; an action like that would have, to quote Mr Albanese,"no excuse".
It wouldn't help the cause of securing a lasting ceasefire, nor agitating towards a free Palestine.
It would be almost universally reviled as an intentional stunt -- just as we ' ve seen it is now that it ' s being @ @ @ @ @ @ @ @ @ @ exploited by the political opponents of a free Palestine as an excuse to shut pro-Palestinian activism down, as Liberal MP Mr Southwick is desperately hoping might now happen.
It ' s a bad idea.
Here ' s some context, then.
Israeli lobbyists and political groups have held a great many meetings in Victorian venues over the past two months, usually for the purpose of lobbying senior Victorian politicians to secure ongoing consent for the indiscriminate assault on Gaza, conducted largely unchecked, after the brutality of 7 October.
Western political consent is very important to the state of Israel.
So they constantly lobby those in power for more of it.
I know about these events because I admit I occasionally receive clandestine information about them.
Similarly, I understand that the activists were acting on information that the Embassy of Israel was hosting an event at the hotel.
The details of the event weren't known to them.
Only who was hosting.
I am led to believe there were, in fact, embassy delegates there @ @ @ @ @ @ @ @ @ @ also happened to be families at the hotel, either attending another booking or simply staying there.
Why were they hoping to protest an Israeli embassy event?
I asked.
They said :"As always, our protest was directed toward those who hold power -- the Israeli officials playing politics with Palestinian lives."They held their rally at about 10 PM to an audience of barely anybody.
We ' re told that families of hostages, returning to the hotel after an event elsewhere, felt unsafe and retreated to a nearby police station.
That ' s terrible.
Now we know it was a mistake that this happened and certainly not a case of any activist intentionally attempting to"organise a demonstration against grieving families", to quote our Prime Minister.
So you are welcome to find it tasteless, but you can also put to bed any discussion of suppressing peaceful protest.
Correct?
Because now we ' re only talking about an incident with bad optics that was a mistake -- an @ @ @ @ @ @ @ @ @ @ It ' s not a case of intentionally targeting families.
Now we know this.
But knowing this doesn't seem to have changed the breathless media coverage on the subject.
The Australian quoted WACA as clarifying what its intentions and planned targets for the protest were, before moving on to quote Executive Council of Australian Jewry CEO Alex Ryvchin competing with others to have the most lurid quote regarding the protesters.
The group of five will meet political leaders during a two-day visit in Canberra.
They are also expected to share their experiences with members of the community in Sydney and Melbourne during a weeklong Australian trip."Our... aim is to get the support of the Australian people and the Australian government... for Israel ' s actions in wiping out Hamas and in our military actions right now to eradicate Hamas."So these families of hostages function as props for lobbyists.
They are being shipped out from Israel in order to curry favour with Australian politicians for more Israeli military action.
And it @ @ @ @ @ @ @ @ @ @ decided that ' s a story as worthy of discussion as -- seemingly deliberately -- misinterpreting this case of crossed wires as an intentional action targeting hostages ' families.
I only wish Australians valued the right to peacefully protest over the ability to engage in backroom political lobbying.
Tom Tanuki is a writer, satirist and anti-fascist activist.
Tom does weekly videos on YouTube commenting on the Australian political fringe.
You can follow Tom on Twitter @tomtanuki.

% matched lemmas: 
\end{textsample}
